\begin{table}[t]
\centering
\caption{Two level-19 carriers, against the observation}
\label{tab:level19_direct}
\footnotesize
\setlength{\tabcolsep}{4pt}
\begin{tabular}{@{}lrrrrrrr@{}}
\toprule
 & & Level 20 & \multicolumn{2}{c}{Direct} & \multicolumn{2}{c}{Aggregate} & Carrier \\
\cmidrule(lr){4-5}\cmidrule(lr){6-7}
Read & $n$ & floor & floor & $/$L20 & floor & $/$L20 & diff. \\
\midrule
nearest cell & 40 & 4.04e-03 & 8.87e-03 & 2.25 & 8.56e-03 & 2.23 & 0.27 \\
interpolating & 40 & 4.04e-03 & 8.87e-03 & 2.25 & 8.56e-03 & 2.23 & 0.27 \\
\bottomrule
\end{tabular}
\vspace{3pt}
\begin{minipage}{\columnwidth}
\footnotesize
Median over the regions of the RMSE of each carrier against the same Sentinel-2 observation, on the intersection mask of Section~\ref{sec:common_reference}, read at the nearest containing cell of that carrier's own level. \emph{Direct} evaluates the Clough--Tocher interpolant of the same native samples at the level-19 cell centres; \emph{aggregate} averages the four level-20 children, which is the exact cell-average of the stored carrier and the object Section~\ref{sec:level_choice} measures. The last column is the RMS difference between the two level-19 carriers in units of the aggregate's own floor: it says how much of the level-19 error budget is decided by which of the two objects is called the level-19 carrier. Ratios are to the level-20 floor of the same read.
\end{minipage}
\end{table}
